\section*{अध्याय \ref{ch:g6:earth-a-planet} --- पृथ्वी, सौर मंडल का एक ग्रह}

\begin{solution}{exo:g6:earth-a-planet:1}
\underline{बुध}, \underline{शुक्र}, \underline{पृथ्वी},
\underline{मंगल}, बृहस्पति, शनि, अरुण, वरुण। \omterm{def:g6:earth-a-planet:planet-types}{चट्टानी ग्रह} पहले चार
हैं।
\end{solution}

\begin{solution}{exo:g6:earth-a-planet:2}
``\omterm{def:g6:earth-a-planet:planet-types}{चट्टानी ग्रह} छोटा और घना होता है, जिसकी सतह ठोस होती है। \omterm{def:g6:earth-a-planet:planet-types}{विशाल ग्रह} बहुत बड़ा
होता है, कम घनत्व वाला, जो ज़्यादातर गैसों और द्रवों से बना होता है।''
\end{solution}

\begin{solution}{exo:g6:earth-a-planet:3}
शुक्र (12\,104~km, पृथ्वी के 12\,756~km के मुक़ाबले)। सबसे बड़ा
बृहस्पति है (142\,984~km); सबसे छोटा बुध है (4\,879~km)।
\end{solution}

\begin{solution}{exo:g6:earth-a-planet:4}
(क) ग़लत: शुक्र पर कोई द्रव पानी नहीं है; वह बहुत ज़्यादा गरम है। (ख) सही
(औसतन $-65$~$^\circ$C, 15~$^\circ$C के मुक़ाबले)। (ग) ग़लत: सूर्य के सबसे पास के
चार ग्रह \omterm{def:g6:earth-a-planet:planet-types}{चट्टानी ग्रह} हैं; \omterm{def:g6:earth-a-planet:planet-types}{विशाल ग्रह} सबसे दूर के चार हैं। (घ) सही: लगभग
0.7~ग्राम प्रति घन सेंटीमीटर।
\end{solution}

\begin{solution}{exo:g6:earth-a-planet:5}
$4879 \div 12756 \approx 0.382$, यानी 38~\%: बुध पृथ्वी की चौड़ाई के एक-तिहाई से
थोड़ा ज़्यादा चौड़ा है।
\end{solution}

\begin{solution}{exo:g6:earth-a-planet:6}
$120536 \div 12756 \approx 9.4$: लगभग 9 पृथ्वियाँ। शनि एक \omterm{def:g6:earth-a-planet:planet-types}{विशाल ग्रह} है।
\end{solution}

\begin{solution}{exo:g6:earth-a-planet:7}
ग्रह~A एक \omterm{def:g6:earth-a-planet:planet-types}{विशाल ग्रह} है: वह बहुत बड़ा है और उसका घनत्व, 1.1~ग्राम प्रति घन
सेंटीमीटर, कम है, सौर मंडल के \omterm{def:g6:earth-a-planet:planet-types}{विशाल ग्रहों} की तरह
(0.7 और 1.7 के बीच)। ग्रह~B एक \omterm{def:g6:earth-a-planet:planet-types}{चट्टानी ग्रह} है: वह छोटा है और उसका घनत्व, 5.0,
ज़्यादा है, \omterm{def:g6:earth-a-planet:planet-types}{चट्टानी ग्रहों} की तरह (3.9 और
5.5 के बीच)।
\end{solution}

\begin{solution}{exo:g6:earth-a-planet:8}
(क) $464 - 15 = 449$~डिग्री ज़्यादा गरम। (ख) $15 - (-65) = 80$~डिग्री ज़्यादा ठंडा।
\end{solution}

\begin{solution}{exo:g6:earth-a-planet:9}
शुक्र का कार्बन डाइऑक्साइड का बहुत घना वायुमंडल है, जो सूर्य के ताप को एक भारी कंबल
की तरह फँसा लेता है और उसे निकलने नहीं देता। बुध, जो सूर्य के ज़्यादा पास है, पर
लगभग कोई हवा नहीं है, इसलिए वह अपना ताप नहीं रख पाता।
\end{solution}

\begin{solution}{exo:g6:earth-a-planet:10}
$228 - 150 = 78$~मिलियन किलोमीटर ज़्यादा दूर। $78 \div 150 = 0.52$: मंगल की
दूरी पृथ्वी की दूरी से 52~\% ज़्यादा है।
\end{solution}

\begin{solution}{exo:g6:earth-a-planet:11}
अरबों वर्ष पहले मंगल पर \omterm{def:g1:landscapes-around-us:water}{नदियाँ}, \omterm{def:g1:landscapes-around-us:water}{झीलें} और \omterm{def:g5:water-shapes-the-land:flood}{बाढ़ें} थीं: वह ज़्यादा गरम और गीला था।
आज उसका पानी जमा हुआ है, ध्रुवीय टोपियों की बर्फ़ में और \omterm{def:g3:continents-and-oceans:poles}{ध्रुवों} के पास ज़मीन के
नीचे। अब उसकी सतह पर द्रव पानी नहीं टिक सकता, क्योंकि मंगल बहुत ठंडा है (औसतन
$-65$~$^\circ$C) और क्योंकि उसकी हवा इतनी पतली है कि पानी जल्दी जम जाता है या
भाप बन जाता है।
\end{solution}

\begin{solution}{exo:g6:earth-a-planet:12}
(क) $8.9 \div 9.8 \approx 0.908$, यानी 91~\%। (ख) जुड़वाँ जैसा: लगभग
एक-सा आकार (पृथ्वी के व्यास का 95~\%), लगभग एक-सा गुरुत्व का
खिंचाव, और वह भी एक \omterm{def:g6:earth-a-planet:planet-types}{चट्टानी ग्रह} है। जुड़वाँ जैसा नहीं: वह बहुत ज़्यादा गरम है
(464~$^\circ$C, 15~$^\circ$C के मुक़ाबले), उसकी हवा कहीं ज़्यादा घनी है और कार्बन
डाइऑक्साइड से बनी है, अम्ल के \omterm{def:g4:clouds-rain-snow:cloud}{बादलों} के साथ, और उस पर कोई द्रव पानी नहीं है।
\end{solution}

\begin{solution}{pb:g6:earth-a-planet:1}
\textbf{1.} शुक्र (108~मिलियन~km), पृथ्वी (150), मंगल (228)।

\textbf{2.} $150 - 108 = 42$: शुक्र 42~मिलियन किलोमीटर ज़्यादा पास है।
$228 - 150 = 78$: मंगल 78~मिलियन किलोमीटर ज़्यादा दूर है।

\textbf{3.} $6792 \div 12756 \approx 0.532$, यानी 53~\%: मंगल पृथ्वी की
चौड़ाई का लगभग आधा है।

\textbf{4.} तीनों \omterm{def:g6:earth-a-planet:planet-types}{चट्टानी ग्रह} हैं: उनके घनत्व (5.2, 5.5 और 3.9
ग्राम प्रति घन सेंटीमीटर) ज़्यादा हैं, और उनके व्यास \omterm{def:g6:earth-a-planet:planet-types}{विशाल ग्रहों} के व्यासों की
तुलना में छोटे हैं।

\textbf{5.} सबसे घनी हवा शुक्र पर है (92~bar); सबसे पतली मंगल पर
(0.01~bar)।

\textbf{6.} $1 \div 0.01 = 100$: मंगल की हवा पृथ्वी की हवा से सौ गुना
कम दबाती है।

\textbf{7.} उसकी कार्बन डाइऑक्साइड की घनी हवा सूर्य के ताप को फँसा लेती है और उसे
निकलने नहीं देती, एक भारी कंबल की तरह।

\textbf{8.} $92 \div 1 = 92$: शुक्र की ज़मीन पर हवा पृथ्वी की ज़मीन की तुलना में
\textbf{92 गुना} ज़ोर से दबाती है। यान की दीवारें बहुत मज़बूत होनी चाहिए ताकि वह
कुचल न जाए (और उसे गरमी और अम्ल के \omterm{def:g4:clouds-rain-snow:cloud}{बादलों} को भी सहना
होगा)।

\textbf{9.} पृथ्वी का औसत तापमान, 15~$^\circ$C, 0~$^\circ$C और
100~$^\circ$C के बीच है: पानी न जमा है न उबल रहा है, इसलिए वह
द्रव रहता है।

\textbf{10.} शुक्र पर, 464~$^\circ$C पर, जो 100~$^\circ$C से बहुत ऊपर है, पानी
तुरंत उबलकर भाप बन जाएगा। मंगल पर, $-65$~$^\circ$C पर, जो
0~$^\circ$C से बहुत नीचे है, वह जमकर बर्फ़ बन जाएगा (और धीरे-धीरे भाप बन जाएगा,
क्योंकि हवा बहुत पतली है)।

\textbf{11.} बहुत पहले मंगल आज से ज़्यादा गरम और गीला था: उसकी सतह पर द्रव पानी
बह सकता था। उसकी \omterm{def:g5:weather-and-climate:climate}{जलवायु} बहुत बदल गई है: अब वह
एक ठंडा, सूखा \omterm{def:g5:weather-and-climate:desert}{मरुस्थल} है।

\textbf{12.} उदाहरण के लिए: ``प्यारे यानो, यहाँ पृथ्वी पर हमारे पास द्रव पानी है,
\omterm{def:g2:water-on-earth:ocean}{महासागरों}, \omterm{def:g1:landscapes-around-us:water}{नदियों} और \omterm{def:g1:landscapes-around-us:water}{झीलों} में। हमारे तापमान हल्के हैं, औसतन 15~$^\circ$C,
पानी के जमने और उबलने के बीच। हम सूर्य से सही दूरी पर हैं, न बहुत पास न बहुत
दूर। और हमारा वायुमंडल हमें गरम रखता है, जीवों को ज़रूरी गैसें देता है, और
पानी को द्रव रहने देता है। इसीलिए यहाँ जीवन संभव है।''
\end{solution}
